% --- Archon physics formalization source begin ---
% archon:physics
% archon:covers PhyXMiniProblems/problem_phyx_mini_0090.lean
% archon:source-report reports/phyx_mini/problem_phyx_mini_0090.source.json

\chapter{Physics problem phyx\_mini\_0090}
\label{ch:PhyXMiniProblems_problem_phyx_mini_0090}

\paragraph{Problem source.}
\#\# Physical scenario

In figure, two isotropic point sources \textbackslash{}( S\_1 \textbackslash{}) and \textbackslash{}( S\_2 \textbackslash{}) emit light in phase at wavelength \textbackslash{}( \textbackslash{}lambda \textbackslash{}) and at the same amplitude. The sources are separated by distance \textbackslash{}( d = 6.00\textbackslash{}lambda \textbackslash{}) on an \textbackslash{}( x \textbackslash{}) axis. A viewing screen is at distance \textbackslash{}( D = 20.0\textbackslash{}lambda \textbackslash{}) from \textbackslash{}( S\_2 \textbackslash{}) and parallel to the \textbackslash{}( y \textbackslash{}) axis. The figure shows two rays reaching point \textbackslash{}( P \textbackslash{}) on the screen, at height \textbackslash{}( y\_p \textbackslash{}).

\#\# Question

What is the phase difference when \textbackslash{}( y\_p = d \textbackslash{})?

\#\# Answer choices

- A: A: 4.80λ

- B: B: 6.00λ

- C: C: 5.80λ

- D: D: 5.00λ

\#\# Figure caption (auxiliary; use the image as primary evidence)

The image depicts a diagram related to wave interference, likely illustrating the double-slit experiment. Here's a detailed breakdown:

- **Axis**: There is a horizontal \textbackslash{}(x\textbackslash{})-axis and a vertical \textbackslash{}(y\textbackslash{})-axis. A screen is positioned parallel to the \textbackslash{}(y\textbackslash{})-axis.

- **Points**: Two points, \textbackslash{}(S\_1\textbackslash{}) and \textbackslash{}(S\_2\textbackslash{}), are marked on the horizontal line, indicating sources. These are represented by red circles.

- **Screen**: A vertical line labeled "Screen" is on the right side, representing where interference patterns might be observed.

- **Point \textbackslash{}(P\textbackslash{})**: A point, \textbackslash{}(P\textbackslash{}), is marked on the screen showing where waves from \textbackslash{}(S\_1\textbackslash{}) and \textbackslash{}(S\_2\textbackslash{}) converge.

- **Lines and Arrows**: Two arrows extend from \textbackslash{}(S\_1\textbackslash{}) and \textbackslash{}(S\_2\textbackslash{}) to point \textbackslash{}(P\textbackslash{}) on the screen, showing the path of waves.

- **Distances**: 
  - The distance between \textbackslash{}(S\_1\textbackslash{}) and \textbackslash{}(S\_2\textbackslash{}) is labeled \textbackslash{}(d\textbackslash{}).
  - The distance from \textbackslash{}(S\_2\textbackslash{}) to the screen along the \textbackslash{}(x\textbackslash{})-axis is labeled \textbackslash{}(D\textbackslash{}).

This figure represents the setup to demonstrate wave interference, showing how waves from two sources meet at a point on a screen.

\#\# Dataset metadata

Wave Optics; Physical Model Grounding Reasoning, Spatial Relation Reasoning

\paragraph{Recorded answer/context.}
C: C: 5.80λ

\paragraph{Figure/image path.}
/root/proposal\_for\_physic/science-mango/phyx\_mini\_run/phyx\_data/test\_image/90.png

\paragraph{Formalization target.}
create a compiling Lean file with sorry bodies at `PhyXMiniProblems/problem\_phyx\_mini\_0090.lean`. The Lean declarations must preserve the physical quantities, dimensions or dimensional roles, figure labels, governing-law hypotheses, and final relation expressed by this problem.
Use Mathlib/Physlib names found through LeanExplore where available. If a physics API is missing, introduce faithful local abstractions rather than scalar placeholder aliases.

\begin{theorem}[Physics formalization target]
\label{thm:physics:phyx_mini_0090:target}
\lean{PhyXMiniProblems.ProblemPhyXMini0090.problem_phyx_mini_0090}
\uses{def:physics:phyx-mini-0090:phyxminiproblems-problemphyxmini0090-twosourceinterferencesetup, def:physics:phyx-mini-0090:phyxminiproblems-problemphyxmini0090-emitsinphaseatsameamplitudeisotropically, def:physics:phyx-mini-0090:phyxminiproblems-problemphyxmini0090-hasdepictedlayout, def:physics:phyx-mini-0090:phyxminiproblems-problemphyxmini0090-matchesproblemreadouts, def:physics:phyx-mini-0090:phyxminiproblems-problemphyxmini0090-haspositivephysicallengths, def:physics:phyx-mini-0090:phyxminiproblems-problemphyxmini0090-satisfieseuclideanraygeometry, def:physics:phyx-mini-0090:phyxminiproblems-problemphyxmini0090-satisfiesmonochromaticpropagationlaw, def:physics:phyx-mini-0090:phyxminiproblems-problemphyxmini0090-phasedifferenceatp, def:physics:phyx-mini-0090:phyxminiproblems-problemphyxmini0090-isnearestdisplayedpathdifference}
The assigned autoformalize agent should translate this physics problem into Lean declarations in the covered file, with theorem and lemma proof bodies written as `by sorry`.
\end{theorem}
\begin{proof}
This is an autoformalization task, not a proof task. Produce statements that can later be proved by the physics prover without weakening the physical model.
\end{proof}

% --- Archon declaration topology begin ---
\section{Declaration topology}
\begin{definition}[Length Quantity]
  \label{def:physics:phyx-mini-0090:phyxminiproblems-problemphyxmini0090-lengthquantity}
  \lean{PhyXMiniProblems.ProblemPhyXMini0090.LengthQuantity}
  A physical quantity carrying the dimension of length
\end{definition}
\begin{proof}
  This is the declared model definition of Length Quantity.
\end{proof}
\begin{definition}[length Readout]
  \label{def:physics:phyx-mini-0090:phyxminiproblems-problemphyxmini0090-lengthreadout}
  \lean{PhyXMiniProblems.ProblemPhyXMini0090.lengthReadout}
  \uses{def:physics:phyx-mini-0090:phyxminiproblems-problemphyxmini0090-lengthquantity}
  The scalar readout of a dimensionful length in a selected length unit
\end{definition}
\begin{proof}
  This is the declared model definition of length Readout.
\end{proof}
\begin{definition}[Source Label]
  \label{def:physics:phyx-mini-0090:phyxminiproblems-problemphyxmini0090-sourcelabel}
  \lean{PhyXMiniProblems.ProblemPhyXMini0090.SourceLabel}
  The two point-source labels shown in the figure
\end{definition}
\begin{proof}
  This is the declared model definition of Source Label.
\end{proof}
\begin{definition}[Radiation Pattern]
  \label{def:physics:phyx-mini-0090:phyxminiproblems-problemphyxmini0090-radiationpattern}
  \lean{PhyXMiniProblems.ProblemPhyXMini0090.RadiationPattern}
  The angular radiation pattern of a source
\end{definition}
\begin{proof}
  This is the declared model definition of Radiation Pattern.
\end{proof}
\begin{definition}[Screen Orientation]
  \label{def:physics:phyx-mini-0090:phyxminiproblems-problemphyxmini0090-screenorientation}
  \lean{PhyXMiniProblems.ProblemPhyXMini0090.ScreenOrientation}
  The qualitative orientation of the viewing screen in the diagram
\end{definition}
\begin{proof}
  This is the declared model definition of Screen Orientation.
\end{proof}
\begin{definition}[Point Light Source]
  \label{def:physics:phyx-mini-0090:phyxminiproblems-problemphyxmini0090-pointlightsource}
  \lean{PhyXMiniProblems.ProblemPhyXMini0090.PointLightSource}
  \uses{def:physics:phyx-mini-0090:phyxminiproblems-problemphyxmini0090-lengthquantity, def:physics:phyx-mini-0090:phyxminiproblems-problemphyxmini0090-radiationpattern}
  A monochromatic point source. Amplitude is intentionally an abstract type: the problem specifies equal amplitudes but supplies no amplitude calibration or unit in which a scalar readout should be taken
\end{definition}
\begin{proof}
  This is the declared model definition of Point Light Source.
\end{proof}
\begin{definition}[Two Source Interference Setup]
  \label{def:physics:phyx-mini-0090:phyxminiproblems-problemphyxmini0090-twosourceinterferencesetup}
  \lean{PhyXMiniProblems.ProblemPhyXMini0090.TwoSourceInterferenceSetup}
  \uses{def:physics:phyx-mini-0090:phyxminiproblems-problemphyxmini0090-lengthquantity, def:physics:phyx-mini-0090:phyxminiproblems-problemphyxmini0090-sourcelabel, def:physics:phyx-mini-0090:phyxminiproblems-problemphyxmini0090-screenorientation, def:physics:phyx-mini-0090:phyxminiproblems-problemphyxmini0090-pointlightsource}
  The physical objects and figure-derived quantities for the two-source setup. screenXCoordinate is a scalar coordinate readout in coordinateUnit, while sourceSeparation, screenDistanceFromS2, observationHeight, wavelengths, and ray lengths are genuine dimensionful quantities
\end{definition}
\begin{proof}
  This is the declared model definition of Two Source Interference Setup.
\end{proof}
\begin{definition}[common Wavelength]
  \label{def:physics:phyx-mini-0090:phyxminiproblems-problemphyxmini0090-commonwavelength}
  \lean{PhyXMiniProblems.ProblemPhyXMini0090.commonWavelength}
  \uses{def:physics:phyx-mini-0090:phyxminiproblems-problemphyxmini0090-lengthquantity, def:physics:phyx-mini-0090:phyxminiproblems-problemphyxmini0090-twosourceinterferencesetup}
  The common wavelength, named using the S₁ source
\end{definition}
\begin{proof}
  This is the declared model definition of common Wavelength.
\end{proof}
\begin{definition}[Emits In Phase At Same Amplitude Isotropically]
  \label{def:physics:phyx-mini-0090:phyxminiproblems-problemphyxmini0090-emitsinphaseatsameamplitudeisotropically}
  \lean{PhyXMiniProblems.ProblemPhyXMini0090.EmitsInPhaseAtSameAmplitudeIsotropically}
  \uses{def:physics:phyx-mini-0090:phyxminiproblems-problemphyxmini0090-twosourceinterferencesetup}
  The source conditions stated in the problem: both sources are isotropic, have the same amplitude and wavelength, and emit in phase. This says nothing about their phase difference after propagation to P
\end{definition}
\begin{proof}
  This is the declared model definition of Emits In Phase At Same Amplitude Isotropically.
\end{proof}
\begin{definition}[Has Depicted Layout]
  \label{def:physics:phyx-mini-0090:phyxminiproblems-problemphyxmini0090-hasdepictedlayout}
  \lean{PhyXMiniProblems.ProblemPhyXMini0090.HasDepictedLayout}
  \uses{def:physics:phyx-mini-0090:phyxminiproblems-problemphyxmini0090-lengthreadout, def:physics:phyx-mini-0090:phyxminiproblems-problemphyxmini0090-twosourceinterferencesetup}
  The categorical and coordinate layout read from the primary image. S₁ is the origin, S₂ lies d to its right, the screen is the vertical line at x = d + D, and P lies on that screen at height yₚ
\end{definition}
\begin{proof}
  This is the declared model definition of Has Depicted Layout.
\end{proof}
\begin{definition}[Matches Problem Readouts]
  \label{def:physics:phyx-mini-0090:phyxminiproblems-problemphyxmini0090-matchesproblemreadouts}
  \lean{PhyXMiniProblems.ProblemPhyXMini0090.MatchesProblemReadouts}
  \uses{def:physics:phyx-mini-0090:phyxminiproblems-problemphyxmini0090-lengthreadout, def:physics:phyx-mini-0090:phyxminiproblems-problemphyxmini0090-twosourceinterferencesetup, def:physics:phyx-mini-0090:phyxminiproblems-problemphyxmini0090-commonwavelength}
  The numerical problem data, expressed covariantly in every length unit: d = 6λ, D = 20λ, and the requested observation height is yₚ = d. No ray-path or phase-difference answer occurs here
\end{definition}
\begin{proof}
  This is the declared model definition of Matches Problem Readouts.
\end{proof}
\begin{definition}[Has Positive Physical Lengths]
  \label{def:physics:phyx-mini-0090:phyxminiproblems-problemphyxmini0090-haspositivephysicallengths}
  \lean{PhyXMiniProblems.ProblemPhyXMini0090.HasPositivePhysicalLengths}
  \uses{def:physics:phyx-mini-0090:phyxminiproblems-problemphyxmini0090-lengthreadout, def:physics:phyx-mini-0090:phyxminiproblems-problemphyxmini0090-sourcelabel, def:physics:phyx-mini-0090:phyxminiproblems-problemphyxmini0090-twosourceinterferencesetup, def:physics:phyx-mini-0090:phyxminiproblems-problemphyxmini0090-commonwavelength}
  Positivity conditions for the physical lengths in the setup
\end{definition}
\begin{proof}
  This is the declared model definition of Has Positive Physical Lengths.
\end{proof}
\begin{definition}[Satisfies Euclidean Ray Geometry]
  \label{def:physics:phyx-mini-0090:phyxminiproblems-problemphyxmini0090-satisfieseuclideanraygeometry}
  \lean{PhyXMiniProblems.ProblemPhyXMini0090.SatisfiesEuclideanRayGeometry}
  \uses{def:physics:phyx-mini-0090:phyxminiproblems-problemphyxmini0090-lengthreadout, def:physics:phyx-mini-0090:phyxminiproblems-problemphyxmini0090-sourcelabel, def:physics:phyx-mini-0090:phyxminiproblems-problemphyxmini0090-twosourceinterferencesetup}
  The geometric-optics law that each ray-path length is the Euclidean distance from its point source to P. This is a generic governing relation and does not contain the requested numerical path or phase difference
\end{definition}
\begin{proof}
  This is the declared model definition of Satisfies Euclidean Ray Geometry.
\end{proof}
\begin{definition}[Satisfies Monochromatic Propagation Law]
  \label{def:physics:phyx-mini-0090:phyxminiproblems-problemphyxmini0090-satisfiesmonochromaticpropagationlaw}
  \lean{PhyXMiniProblems.ProblemPhyXMini0090.SatisfiesMonochromaticPropagationLaw}
  \uses{def:physics:phyx-mini-0090:phyxminiproblems-problemphyxmini0090-lengthreadout, def:physics:phyx-mini-0090:phyxminiproblems-problemphyxmini0090-sourcelabel, def:physics:phyx-mini-0090:phyxminiproblems-problemphyxmini0090-twosourceinterferencesetup}
  For monochromatic propagation, a path of length r accumulates phase 2π r / λ, regarded as an element of Real.Angle. The law is stated for each source separately and does not assert the requested phase difference
\end{definition}
\begin{proof}
  This is the declared model definition of Satisfies Monochromatic Propagation Law.
\end{proof}
\begin{definition}[phase At P]
  \label{def:physics:phyx-mini-0090:phyxminiproblems-problemphyxmini0090-phaseatp}
  \lean{PhyXMiniProblems.ProblemPhyXMini0090.phaseAtP}
  \uses{def:physics:phyx-mini-0090:phyxminiproblems-problemphyxmini0090-sourcelabel, def:physics:phyx-mini-0090:phyxminiproblems-problemphyxmini0090-twosourceinterferencesetup}
  The total phase of the wave from one source when it reaches P
\end{definition}
\begin{proof}
  This is the declared model definition of phase At P.
\end{proof}
\begin{definition}[phase Difference At P]
  \label{def:physics:phyx-mini-0090:phyxminiproblems-problemphyxmini0090-phasedifferenceatp}
  \lean{PhyXMiniProblems.ProblemPhyXMini0090.phaseDifferenceAtP}
  \uses{def:physics:phyx-mini-0090:phyxminiproblems-problemphyxmini0090-twosourceinterferencesetup, def:physics:phyx-mini-0090:phyxminiproblems-problemphyxmini0090-phaseatp}
  The directed phase difference phase(S₁ at P) - phase(S₂ at P), modulo 2π. The figure makes the S₁ ray longer, so this orientation agrees with the positive path difference used by the displayed choices
\end{definition}
\begin{proof}
  This is the declared model definition of phase Difference At P.
\end{proof}
\begin{definition}[path Difference In Wavelengths]
  \label{def:physics:phyx-mini-0090:phyxminiproblems-problemphyxmini0090-pathdifferenceinwavelengths}
  \lean{PhyXMiniProblems.ProblemPhyXMini0090.pathDifferenceInWavelengths}
  \uses{def:physics:phyx-mini-0090:phyxminiproblems-problemphyxmini0090-lengthreadout, def:physics:phyx-mini-0090:phyxminiproblems-problemphyxmini0090-twosourceinterferencesetup, def:physics:phyx-mini-0090:phyxminiproblems-problemphyxmini0090-commonwavelength}
  The geometric path difference (r₁-r₂)/λ, in wavelength units
\end{definition}
\begin{proof}
  This is the declared model definition of path Difference In Wavelengths.
\end{proof}
\begin{definition}[Answer Choice]
  \label{def:physics:phyx-mini-0090:phyxminiproblems-problemphyxmini0090-answerchoice}
  \lean{PhyXMiniProblems.ProblemPhyXMini0090.AnswerChoice}
  Labels of the four multiple-choice answers
\end{definition}
\begin{proof}
  This is the declared model definition of Answer Choice.
\end{proof}
\begin{definition}[answer Path Difference In Wavelengths]
  \label{def:physics:phyx-mini-0090:phyxminiproblems-problemphyxmini0090-answerpathdifferenceinwavelengths}
  \lean{PhyXMiniProblems.ProblemPhyXMini0090.answerPathDifferenceInWavelengths}
  \uses{def:physics:phyx-mini-0090:phyxminiproblems-problemphyxmini0090-answerchoice}
  The path difference printed beside each answer, measured in wavelengths. These are the answer-table readouts, not assumptions about the physical rays
\end{definition}
\begin{proof}
  This is the declared model definition of answer Path Difference In Wavelengths.
\end{proof}
\begin{definition}[Is Nearest Displayed Path Difference]
  \label{def:physics:phyx-mini-0090:phyxminiproblems-problemphyxmini0090-isnearestdisplayedpathdifference}
  \lean{PhyXMiniProblems.ProblemPhyXMini0090.IsNearestDisplayedPathDifference}
  \uses{def:physics:phyx-mini-0090:phyxminiproblems-problemphyxmini0090-twosourceinterferencesetup, def:physics:phyx-mini-0090:phyxminiproblems-problemphyxmini0090-pathdifferenceinwavelengths, def:physics:phyx-mini-0090:phyxminiproblems-problemphyxmini0090-answerchoice, def:physics:phyx-mini-0090:phyxminiproblems-problemphyxmini0090-answerpathdifferenceinwavelengths}
  A displayed choice is selected when its wavelength count is at least as close to the exact geometric path difference as every other displayed value. This models the rounding implicit in the recorded answer 5.80λ
\end{definition}
\begin{proof}
  This is the declared model definition of Is Nearest Displayed Path Difference.
\end{proof}
\begin{lemma}[normalized Ray Path Lengths]
  \label{lem:physics:phyx-mini-0090:phyxminiproblems-problemphyxmini0090-normalizedraypathlengths}
  \lean{PhyXMiniProblems.ProblemPhyXMini0090.normalizedRayPathLengths}
  \uses{def:physics:phyx-mini-0090:phyxminiproblems-problemphyxmini0090-lengthreadout, def:physics:phyx-mini-0090:phyxminiproblems-problemphyxmini0090-twosourceinterferencesetup, def:physics:phyx-mini-0090:phyxminiproblems-problemphyxmini0090-commonwavelength, def:physics:phyx-mini-0090:phyxminiproblems-problemphyxmini0090-hasdepictedlayout, def:physics:phyx-mini-0090:phyxminiproblems-problemphyxmini0090-matchesproblemreadouts, def:physics:phyx-mini-0090:phyxminiproblems-problemphyxmini0090-haspositivephysicallengths, def:physics:phyx-mini-0090:phyxminiproblems-problemphyxmini0090-satisfieseuclideanraygeometry}
  At d = yₚ = 6λ and D = 20λ, the two ray lengths are respectively √712 λ and √436 λ. This is derived from the depicted coordinates and the Euclidean ray-length law
\end{lemma}
\begin{proof}
  The conclusion follows from the governing relations and figure readouts named in the statement of normalized Ray Path Lengths.
\end{proof}
\begin{lemma}[path Difference In Wavelengths eq]
  \label{lem:physics:phyx-mini-0090:phyxminiproblems-problemphyxmini0090-pathdifferenceinwavelengths-eq}
  \lean{PhyXMiniProblems.ProblemPhyXMini0090.pathDifferenceInWavelengths_eq}
  \uses{def:physics:phyx-mini-0090:phyxminiproblems-problemphyxmini0090-twosourceinterferencesetup, def:physics:phyx-mini-0090:phyxminiproblems-problemphyxmini0090-hasdepictedlayout, def:physics:phyx-mini-0090:phyxminiproblems-problemphyxmini0090-matchesproblemreadouts, def:physics:phyx-mini-0090:phyxminiproblems-problemphyxmini0090-haspositivephysicallengths, def:physics:phyx-mini-0090:phyxminiproblems-problemphyxmini0090-satisfieseuclideanraygeometry, def:physics:phyx-mini-0090:phyxminiproblems-problemphyxmini0090-pathdifferenceinwavelengths}
  The exact geometric path difference in wavelength units
\end{lemma}
\begin{proof}
  The conclusion follows from the governing relations and figure readouts named in the statement of path Difference In Wavelengths eq.
\end{proof}
\begin{lemma}[phase Difference At P eq]
  \label{lem:physics:phyx-mini-0090:phyxminiproblems-problemphyxmini0090-phasedifferenceatp-eq}
  \lean{PhyXMiniProblems.ProblemPhyXMini0090.phaseDifferenceAtP_eq}
  \uses{def:physics:phyx-mini-0090:phyxminiproblems-problemphyxmini0090-twosourceinterferencesetup, def:physics:phyx-mini-0090:phyxminiproblems-problemphyxmini0090-emitsinphaseatsameamplitudeisotropically, def:physics:phyx-mini-0090:phyxminiproblems-problemphyxmini0090-hasdepictedlayout, def:physics:phyx-mini-0090:phyxminiproblems-problemphyxmini0090-matchesproblemreadouts, def:physics:phyx-mini-0090:phyxminiproblems-problemphyxmini0090-haspositivephysicallengths, def:physics:phyx-mini-0090:phyxminiproblems-problemphyxmini0090-satisfieseuclideanraygeometry, def:physics:phyx-mini-0090:phyxminiproblems-problemphyxmini0090-satisfiesmonochromaticpropagationlaw, def:physics:phyx-mini-0090:phyxminiproblems-problemphyxmini0090-phasedifferenceatp}
  Equal source phases and the monochromatic propagation law turn the geometric path difference into the corresponding phase difference modulo 2π
\end{lemma}
\begin{proof}
  The conclusion follows from the governing relations and figure readouts named in the statement of phase Difference At P eq.
\end{proof}
% --- Archon declaration topology end ---
% --- Archon physics formalization source end ---
